\cotefile{10}
\scanpage{69}
% archive sheet number: 68

$E$ sur $K$, munis d'un \deleted{\unsure{endomor}} \uline{automorphisme} $\handwritten{\sigma}$-linéaire $F_E$\handwritten{;} \deleted{\unsure{$E\quad E$}} qu'on a intérêt à interpréter comme un isomorphisme $K$-linéaire
\[
F_E : E^{\handwritten{\sigma}} = E\otimes_K(K,\handwritten{\sigma})\handwritten{\longrightarrow} E\ .
\]
Les opérations tensorielles sont définies par transport de structure (cf. cas d'un groupe, ici $\underline{Z}$, opérant sur un corps ....). Les F-isocristaux effectifs s'interprètent alors comme les $(E,F_E)$ tels que $F_E$ laisse invariant un réseau, ce qui signifie aussi que l'ensemble des itérés $F_E^n$, $n\handwritten{\geq}0$, est \uline{borné} dans \unsure{l'e|le}\deleted{\unsure{x}} vectoriel $\operatorname{End}(E)$. Tout$(E,F_E\unsure{)}$ devient \inserted{évidemment} effectif après multiplication de $F_E$ par une puissance convenable de $p$, ce qui est une autre façon de dire que la catégorie des F-iso\unsure{c|w}ristaux se déduit de celle des F-\deleted{\unsure{si}}\handwritten{\unsure{is|si}}ocristaux effectifs \deleted{par} en rendant \unsure{inversible|ènversible} $K\handwritten{(}-1)$ .

\handwritten{4.} La catégorie tannakienne $\mathrm{Fcriso}(k)$ est un invariant arith\inserted{\handwritten{m}}étique intéressant attaché à $k$ (fonctoriellement); sa \deleted{donnée} connaissance équivaut à celle de la gerbe associée (sur $\underline{Q}_p$), soit $\underline{G}(k)$\handwritten{\unsure{.|;}} Parmi les structures supplémentaires qu'il convient d'étudier en même temps que cette gerbe, il faut inclure la structure d'effectivité de la catégorie $\mathrm{Fc}\handwritten{\mathrm{ri}}\mathrm{so}(k)$ de\deleted{s} ses représentations linéaires, celle de l\deleted{\unsure{à|a}}\handwritten{'}objet de \deleted{\unsure{F}}\handwritten{T}ate inverse (définissant un homomorphisme canonique de $\underline{G}(k)$ vers la gerbe triviale de groupe $\underline{G}_m$), enfin pour $k$ parfait tout au moins, la donnée du foncteur fibre naturel sur $K(k)$ (qui est donc une section de la gerbe $\underline{G}(k)$ sur $K(k)$); pour $k$ non parfait, il con\deleted{\unsure{nai}}viendrait \deleted{de} d'étudier de même la section de $\underline{\unsure{G|E}}(k)$ sur le corps des fractions $K$ de n'importe quel $p$-anneau $W$ de corps résiduel $k$.

L'intérêt de la notion de F-cristal \deleted{non dégénéré} provient du fait que l'on \deleted{\unsure{n}}\unsure{a|â} un foncteur ``cohomologie cristalline''
\[
X\handwritten{\unsure{\longmapsto|\mapstochar\relbar}} H^{\handwritten{*}\deleted{\unsure{k|i}}}_{\mathrm{cris}}(X)
\]
